\sectiontheme{goblue}
\subsection{Create \& Update Actions}
\label{sec:go-actions}

Every entity gets two generated functions --- \texttt{\{Entity\}CreateFn}
and \texttt{\{Entity\}UpdateFn} --- plus a small actions bundle wiring
them up as the defaults, mirroring the pattern
\href{https://github.com/torabian/fireback}{fireback} uses.

\subsubsection{Request body formats}
Every generated Go action --- Create/Update included --- reads its body
through \texttt{emigo.BindGinRequestBody}, which picks the decoder from
the request's own \texttt{Content-Type} rather than assuming JSON:
\begin{itemize}[leftmargin=1.1em,itemsep=2pt,topsep=2pt]
  \item \texttt{application/json} --- the default, and what every
    generated client sends.
  \item \texttt{application/yaml}/\texttt{x-yaml}/\texttt{text/yaml}
    --- the same struct, from YAML.
  \item \texttt{application/xml}/\texttt{text/xml}.
  \item \texttt{multipart/form-data} --- each non-file field round-trips
    through JSON into the struct; each uploaded file becomes a
    \texttt{BindMultipartFile} (\texttt{filename}/\texttt{mime}/
    \texttt{data}), unless the host app registers its own
    \texttt{ConvertMultipartFile} (fireback, for instance, turns one into
    an \texttt{XFile}).
  \item \texttt{application/x-www-form-urlencoded} --- a plain HTML
    \texttt{<form method="post">} postback works unmodified, no
    JavaScript or fetch call required: the browser's own form submission
    already sends this content type, and every field lands in the same
    generated struct a JSON body would.
\end{itemize}
Only \texttt{POST}/\texttt{PATCH}/\texttt{PUT} read a body at all --- a
no-op for every other method, matching plain HTTP. The response side
mirrors this: \texttt{RenderGinResult} content-negotiates off the
request's \texttt{Accept} header (JSON, YAML, TOML, or CSV; JSON when
unset or unrecognized), so a caller gets the format it asked for without
the handler itself knowing or caring which one that is.

\subsubsection{The update input: optional, all the way down}
A struct-based \texttt{gorm.Updates()} call silently drops zero-valued
fields, so it can't tell ``the caller didn't mention this field'' apart
from ``the caller explicitly set it to zero.'' Every entity therefore
gets a second generated type, \texttt{\{Entity\}UpdateInput}, where
\emph{every} field --- even a plain, always-present scalar on the entity
itself --- is wrapped optional:
\begin{lstlisting}
type Entity1EntityUpdateInput struct {
  Title  emigo.Nullable[string] `json:"title"`
  Items  emigo.ArrayNullable[Entity1EntityItems] `json:"items"`
  Owner  emigo.OneNullable[Entity2Entity]         `json:"owner"`
  Complex1 *Money                                 `json:"complex1"`
}
\end{lstlisting}
\texttt{id}/\texttt{uniqueId} are left out entirely --- immutable
identity, not something an update patches. A \texttt{complex} field (no
\texttt{?} counterpart exists) becomes a plain pointer instead;
\texttt{nil} means untouched. Critically, \texttt{array}/
\texttt{collection}/\texttt{one} fields reuse the \emph{real} row/target
type directly, not a second, separately-optional clone --- only whether
the field itself was touched (\texttt{IsSet()}) is optional; an item
included in a replace/append batch still has to be given in full, exactly
like Create, because the reconcile helpers persist each item with a plain
\texttt{Save()} that writes every field.

\subsubsection{The reconcile helpers (\texttt{emigorm})}
Gorm's own \texttt{Association().Replace()} was verified \textbf{unreliable}
for a has-many \texttt{array} relation: an incoming item whose primary key
already exists silently keeps its \emph{old} content, and items dropped
from the list are never deleted --- just orphaned. So
\texttt{github.com/torabian/emi/emigorm} (kept out of the wasm-safe
\texttt{emigo} core since it depends on \texttt{gorm.io/gorm}) provides
three reconcile functions instead:

\begin{sidebar}{One reconciler per relation kind}
\textbf{ReconcileHasMany} (\texttt{array}) --- diffs existing children
against the incoming list by \texttt{UniqueId}, hard-deletes any row
missing from a \texttt{"replace"} batch, \texttt{Save()}s every incoming
item. \texttt{"append"} skips the diff/delete step.

\textbf{ReconcileManyToMany} (\texttt{collection}) --- gorm's Association
API \emph{is} reliable here, so this upserts inline target values first,
then delegates to \texttt{Association().Replace()}/\texttt{.Append()}.

\textbf{ReconcileOne} (\texttt{one}) --- \texttt{"select"} resolves an
existing target's \texttt{Id} from its \texttt{UniqueId} and leaves it
untouched; anything else upserts the inline value.
\end{sidebar}

All three resolve an incoming item's real \texttt{Id} by matching
\texttt{UniqueId} against an existing row \emph{before} calling
\texttt{Save()} --- otherwise every ``update'' would silently insert a
duplicate.

\subsubsection{Testing}
Every relation is verified two ways: \textbf{structurally}, with no
database at all, via \texttt{gorm.io/gorm/schema.Parse} (confirms every
\texttt{has\_many}/\texttt{many\_to\_many}/\texttt{belongs\_to} resolves
correctly in milliseconds); and \textbf{end to end against real
Postgres} (\texttt{examples/emi-entity/sdk/*\_test.go}) --- cascading
create, partial update, array replace-with-orphan-deletion, collection
append, and one re-select are all covered. Since \texttt{UniqueId}'s
column default is Postgres-specific, the live tests need a real Postgres
instance (\texttt{make entity-db-up}); MySQL is wired up too but expected
to fail on the same \texttt{gen\_random\_uuid()} default --- an accepted
tradeoff, not something actively kept green.
